\documentclass[12pt, a4paper]{article}

% ── Packages ─────────────────────────────────────────────────────────────────
\usepackage[T1]{fontenc}
\usepackage[utf8]{inputenc}
\usepackage{lmodern}
\usepackage{microtype}
\usepackage[top=2.5cm, bottom=2.5cm, left=3cm, right=3cm]{geometry}
\setlength{\headheight}{13.6pt}
\addtolength{\topmargin}{-1.6pt}
\usepackage{amsmath, amssymb}
\usepackage{booktabs}
\usepackage{array}
\usepackage{multirow}
\usepackage{graphicx}
\usepackage{caption}
\usepackage{subcaption}
\usepackage{fancyhdr}
\usepackage{parskip}
\usepackage{enumitem}
\usepackage{titlesec}
\usepackage{listings}
\usepackage{makecell}
\usepackage[hidelinks]{hyperref}

% ── Section formatting ────────────────────────────────────────────────────────
\titleformat{\section}{\large\bfseries}{\thesection.}{0.5em}{}[\titlerule]
\titleformat{\subsection}{\normalsize\bfseries}{\thesubsection.}{0.5em}{}
\titleformat{\subsubsection}{\normalsize\itshape}{\thesubsubsection.}{0.5em}{}

% ── Header / Footer ───────────────────────────────────────────────────────────
\pagestyle{fancy}
\fancyhf{}
\rhead{\small Korean Food Freshness Classification}
\lhead{\small Preliminary Report --- 2026}
\cfoot{\thepage}
\renewcommand{\headrulewidth}{0.4pt}

% ── Graphics path ─────────────────────────────────────────────────────────────
\graphicspath{{figures/}}

% ── Code style ────────────────────────────────────────────────────────────────
\lstset{
    basicstyle=\small\ttfamily,
    frame=single,
    language=Python,
    breaklines=true,
    showstringspaces=false,
    numbers=none,
}

% ── Caption setup ─────────────────────────────────────────────────────────────
\captionsetup{font=small, labelfont=bf, skip=6pt}

% ── Table column types ────────────────────────────────────────────────────────
\newcolumntype{L}[1]{>{\raggedright\arraybackslash}p{#1}}
\newcolumntype{C}[1]{>{\centering\arraybackslash}p{#1}}

% ─────────────────────────────────────────────────────────────────────────────
\begin{document}

% ── Title Page ────────────────────────────────────────────────────────────────
\begin{titlepage}
\centering
\vspace*{2.5cm}

{\LARGE\bfseries Freshness Classification of Korean Food}\\[3.5cm]

{\normalsize Digital Image Processing}\\[0.4em]
{\normalsize 2026}\\[3cm]

\hrule
\end{titlepage}

\tableofcontents
\newpage

% ─────────────────────────────────────────────────────────────────────────────
\section{Datasets}
% ─────────────────────────────────────────────────────────────────────────────

\subsection{Available Datasets}

Seven publicly available datasets were identified and surveyed.
Table~\ref{tab:datasets} summarizes their properties.
Together they cover fruits, vegetables, and citrus varieties, with fresh/rotten
binary labels and, in one case, a three-class annotation including formalin-mixed
samples.
One dataset (Ripening Stages) uses an object detection format (YOLOv5) rather
than image classification and is therefore excluded from training experiments.

\begin{table}[h]
\centering
\caption{Datasets surveyed. $^{\dagger}$Object detection format --- excluded from classification experiments.}
\label{tab:datasets}
\begin{tabular}{L{3.0cm} L{3.6cm} C{2.6cm} C{1.4cm} C{1.5cm}}
\toprule
\textbf{Dataset} & \textbf{Categories} & \textbf{Classes} & \textbf{Images} & \textbf{Source} \\
\midrule
FruitVision
    & Apple, Banana, Grape, Mango, Orange
    & Fresh / Formalin / Rotten
    & 10,154 & Mendeley \\[4pt]
Lemon Varieties
    & 7 lemon varieties
    & Fresh / Rotten
    & 1,956 & Mendeley \\[4pt]
Ripening Stages$^{\dagger}$
    & Strawberry, Avocado
    & Unripe / Partial / Ripe / Rotten
    & --- & Mendeley \\[4pt]
Fruits Classification
    & Peach, Pomegranate, Strawberry
    & Fresh / Rotten
    & 1,655 & Mendeley \\[4pt]
Fruits Quality
    & 12 mixed fruits
    & Fresh / Rotten
    & 359 & Kaggle \\[4pt]
Fruits Fresh/Rotten
    & Apple, Banana, Orange
    & Fresh / Rotten
    & 13,599 & Kaggle \\[4pt]
Fresh \& Stale
    & Apple, Banana, Cucumber, Okra, Orange, Potato, Tomato
    & Fresh / Rotten
    & 27,317 & Kaggle \\
\midrule
\multicolumn{3}{l}{\textbf{Total (classification datasets)}}
    & \textbf{55,040} & \\
\bottomrule
\end{tabular}
\end{table}

None of the available datasets contains Korean food items.
The FruitVision dataset is the most challenging due to its three-class structure,
where formalin-mixed samples are visually similar to fresh ones --- making it
a useful proxy for detecting subtle spoilage indicators, analogous to the
difficulty expected in fermented Korean food.

\subsection{Korean Food Dataset --- Proposal}

No labeled dataset for Korean food freshness currently exists in public
repositories.
Table~\ref{tab:korean} lists the target food categories and classification
states for the proposed dataset.

\begin{table}[h]
\centering
\caption{Target categories for the proposed Korean food freshness dataset.}
\label{tab:korean}
\begin{tabular}{L{3cm} C{2.5cm} C{3.2cm} L{4cm}}
\toprule
\textbf{Category} & \textbf{Romanization} & \textbf{States} & \textbf{Notes} \\
\midrule
Kimchi           & kimchi       & Fresh / Rotten             & Color and mold during spoilage \\
Rice cake        & tteok        & Fresh / Rotten             & Mold visible on surface \\
Side dishes      & banchan      & Fresh / Rotten             & Vegetable-based, wilting \\
Cooked rice      & bap          & Fresh / Rotten             & Mold indicator \\
Korean apple     & sagwa        & Fresh / Rotten             & Similar to Western apple \\
Mandarin (Jeju)  & gyul         & Fresh / Semifresh / Rotten & Jeju variety, citrus \\
Shine muscat     & shine muscat & Fresh / Semifresh / Rotten & High-value grape variety \\
Strawberry       & ttalgi       & Fresh / Rotten             & Fast spoilage \\
\bottomrule
\end{tabular}
\end{table}

% ─────────────────────────────────────────────────────────────────────────────
\section{Methodology}
% ─────────────────────────────────────────────────────────────────────────────

\subsection{Overall Pipeline}

The classification pipeline consists of four sequential stages:

\begin{center}
\begin{tabular}{clL{7.5cm}}
\toprule
\textbf{Stage} & \textbf{Component} & \textbf{Description} \\
\midrule
1 & Data preparation   & Raw datasets reorganized into unified
                          \texttt{food/state/} folder structure \\
2 & Feature extraction & Pre-trained CNN backbone (ImageNet weights),
                          frozen or partially unfrozen \\
3 & Projection head    & $d_\text{out} \!\to\! 256 \!\to\! 128$
                          linear projection with ReLU \\
4 & Classification     & Linear layer $128 \!\to\! C$
                          with cross-entropy loss \\
\bottomrule
\end{tabular}
\end{center}

All input images are resized to $224 \times 224$ pixels.
Training images undergo data augmentation (random resized crop, horizontal flip,
$\pm 30^{\circ}$ rotation, color jitter on brightness/contrast/saturation/hue,
Gaussian blur) to increase robustness to lighting and viewpoint variation.
Validation images are center-cropped without augmentation.
All images are normalized with ImageNet channel statistics
($\mu = [0.485, 0.456, 0.406]$, $\sigma = [0.229, 0.224, 0.225]$).

\subsection{Backbone Architectures}

Six backbone architectures are considered.
Table~\ref{tab:backbones} summarizes their properties.
Preliminary experiments are conducted with ResNet-18 as the primary backbone.
EfficientNet-B0 is identified as the recommended architecture for the final
Korean food classifier based on its accuracy-to-parameter ratio and demonstrated
suitability for small datasets.

\begin{table}[h]
\centering
\caption{Backbone architectures considered in this study.}
\label{tab:backbones}
\begin{tabular}{lC{2cm}C{2cm}C{2.2cm}C{1.5cm}}
\toprule
\textbf{Model} & \textbf{Params} & \textbf{Out dim} &
\makecell{\textbf{ImageNet}\\\textbf{Top-1}} & \textbf{Family} \\
\midrule
ResNet-18       & 11.2M &  512 & 69.8\% & ResNet      \\
ResNet-34       & 21.3M &  512 & 73.3\% & ResNet      \\
ResNet-50       & 25.6M & 2048 & 76.1\% & ResNet      \\
MobileNetV3-S   &  2.5M &  576 & 67.7\% & MobileNet   \\
EfficientNet-B0 &  5.3M & 1280 & 77.1\% & EfficientNet \\
EfficientNet-B2 &  9.1M & 1408 & 80.1\% & EfficientNet \\
\bottomrule
\end{tabular}
\end{table}

\subsection{Training Conditions}

Four training conditions form an ablation study to determine the optimal degree
of backbone fine-tuning.
Conditions C3 and C4 (with projection head) are the focus of this study,
as the additional non-linear projection consistently improves performance over
a direct linear probe.

\begin{table}[h]
\centering
\caption{Training conditions (ablation study).}
\label{tab:conditions}
\begin{tabular}{llllL{4cm}}
\toprule
\textbf{ID} & \textbf{Name} & \textbf{Backbone} & \textbf{Head} &
\textbf{Description} \\
\midrule
C1 & \texttt{frozen}       & Frozen      & Linear only         & No backbone adaptation \\
C2 & \texttt{layer4}       & Layer4 free & Linear only         & Partial fine-tuning, no projection \\
C3 & \texttt{head\_frozen} & Frozen      & Projection + Linear & Full head, frozen backbone \\
C4 & \texttt{head\_layer4} & Layer4 free & Projection + Linear & Full head + partial fine-tuning \\
\bottomrule
\end{tabular}
\end{table}

All models are trained with Adam optimizer
($\text{lr} = 10^{-3}$, $\text{lr}_\text{backbone} = 10^{-5}$,
weight decay $= 10^{-4}$) and cosine annealing over 50 epochs.
Batch size is 32.

\subsection{Label Modes}

Two classification targets are evaluated:

\begin{itemize}
    \item \textbf{state}: classes correspond to freshness states only
    (fresh / rotten, or fresh / formalin / rotten). The model learns
    state-invariant features across all food categories.
    \item \textbf{fruit\_state}: classes correspond to food-state combinations
    (e.g., \texttt{apple\_fresh}, \texttt{apple\_rotten}). The model
    simultaneously learns food identity and freshness state, enabling a richer
    confusion matrix analysis.
\end{itemize}

\subsection{Evaluation Metrics}

Given the potential class imbalance in real-world Korean food datasets,
accuracy alone is insufficient.
We report:
\begin{itemize}
    \item \textbf{Macro F1-score}: harmonic mean of precision and recall,
    averaged equally across classes regardless of class frequency.
    \item \textbf{Macro Precision and Recall}: sensitivity to false positives
    and false negatives per class.
    \item \textbf{Confusion matrix}: row-normalized per true class, to identify
    systematic misclassification patterns.
\end{itemize}
For food safety applications, recall on the \textit{rotten} class is
particularly critical: a false negative (rotten food classified as fresh)
carries higher risk than a false positive.

% ─────────────────────────────────────────────────────────────────────────────
\section{Preliminary Experiments}
% ─────────────────────────────────────────────────────────────────────────────

\subsection{Setup}

All experiments in this section use ResNet-18 with ImageNet pre-trained weights.
Conditions C3 and C4 (with projection head) are the primary focus.
A stratified 80/20 train/validation split is applied to all datasets.
Experiments marked \textit{pending} are in progress at the time of this report.

\subsection{Results}

Table~\ref{tab:results} reports the best macro F1-score achieved per dataset
and condition combination.

\begin{table}[h]
\centering
\caption{Best macro F1-score (\%) --- ResNet-18, conditions C3 and C4.
\textit{Pending}: experiment in progress.
N/A: fruit\_state not applicable (single-fruit dataset).
Bold: best result per dataset.}
\label{tab:results}
\setlength{\tabcolsep}{6pt}
\begin{tabular}{lC{2cm}C{2.2cm}C{2cm}C{2.2cm}}
\toprule
\multirow{2}{*}{\textbf{Dataset}}
    & \multicolumn{2}{c}{\textbf{C3: head\_frozen}}
    & \multicolumn{2}{c}{\textbf{C4: head\_layer4}} \\
\cmidrule(lr){2-3}\cmidrule(lr){4-5}
    & state & fruit\_state & state & fruit\_state \\
\midrule
FruitVision (3-class)
    & 82.6 & 86.2 & \textbf{88.0} & \textit{pending} \\
Fruits Classification
    & 94.0 & \textit{pending}
    & \textbf{95.8} & \textit{pending} \\
Lemon Varieties
    & \textit{pending} & N/A & \textbf{98.2} & N/A \\
Fruits Fresh/Rotten
    & \textit{pending} & 99.2
    & \textit{pending} & \textbf{99.5} \\
Fresh \& Stale (9 categories)
    & \textit{pending} & \textit{pending}
    & \textit{pending} & \textit{pending} \\
Fruits Quality
    & \textit{pending} & \textit{pending}
    & \textit{pending} & \textit{pending} \\
\bottomrule
\end{tabular}
\end{table}

\subsection{Confusion Matrix Analysis}

Figure~\ref{fig:cm} shows the row-normalized confusion matrices for two
representative completed experiments.

\begin{figure}[h]
\centering
\includegraphics[width=\textwidth]{confusion_matrix.png}
\caption{Row-normalized confusion matrices (final epoch).
The formalin class in FruitVision shows the highest misclassification rate,
predominantly confused with fresh --- analogous to the challenge of detecting
subtle spoilage in Korean fermented foods.
The Kaggle Fruits Fresh/Rotten matrix (6 classes, fruit\_state mode) shows
near-perfect diagonal structure with essentially no cross-fruit confusion.}
\label{fig:cm}
\end{figure}

\subsection{Learning Curves}

Figures~\ref{fig:acc}--\ref{fig:recall} show the per-epoch metrics for all
completed runs, grouped by dataset.
Each curve corresponds to one training condition and backbone combination.

\begin{figure}[h]
\centering
\includegraphics[width=\textwidth]{acc_curves.png}
\caption{Validation accuracy over training epochs per dataset.}
\label{fig:acc}
\end{figure}

\begin{figure}[h]
\centering
\includegraphics[width=\textwidth]{f1_curves.png}
\caption{Macro-averaged F1 score over training epochs per dataset.}
\label{fig:f1}
\end{figure}

\begin{figure}[h]
\centering
\includegraphics[width=\textwidth]{precision_curves.png}
\caption{Macro-averaged precision over training epochs per dataset.}
\label{fig:prec}
\end{figure}

\begin{figure}[h]
\centering
\includegraphics[width=\textwidth]{recall_curves.png}
\caption{Macro-averaged recall over training epochs per dataset.}
\label{fig:recall}
\end{figure}

\subsection{Key Findings}

The following observations are drawn from completed experiments:

\begin{enumerate}

\item \textbf{Fine-tuning is critical for subtle spoilage detection.}
On FruitVision (3-class), the gap between C3 (frozen backbone, 82.6\%)
and C4 (layer4 fine-tuned, 88.0\%) is 5.4 percentage points.
For binary problems, the gap is smaller but consistent:
Fruits Classification achieves 94.0\% (C3) vs.\ 95.8\% (C4).

\item \textbf{Formalin is the hardest class to detect.}
In FruitVision (C4, state mode), per-class F1 at 50 epochs:
rotten 96.4\%, fresh 85.4\%, formalin 82.3\%.
The safety-critical error is formalin predicted as fresh: 135 out of 635
formalin samples (21.3\%) pass undetected.
In fruit\_state mode (15 classes, C3), this dangerous error is reduced:
only 9 out of 129 apple\_formalin samples are predicted as apple\_fresh (7.0\%).
However, fresh apples are frequently misclassified as formalin (69/153 = 45\%),
indicating the frozen backbone struggles to separate the two visually similar
states and errs on the side of over-flagging.

\item \textbf{fruit\_state mode improves FruitVision with a frozen backbone.}
C3 achieves 82.6\% in state mode (3 classes) but 86.2\% in fruit\_state mode
(15 classes). The additional fruit-identity signal helps the frozen backbone
form better decision boundaries even without fine-tuning.

\item \textbf{Transfer learning reaches near-perfect accuracy on clean datasets.}
Fruits Fresh/Rotten (apple, banana, orange) achieves 99.5\% F1 with C4 in
fruit\_state mode (6 classes), with the confusion matrix showing near-zero
off-diagonal entries --- the model jointly learns fruit identity and freshness
state without error.

\item \textbf{Small datasets converge early and require early stopping.}
Lemon Varieties (1,956 images) peaks at epoch 25 (98.2\% F1) and shows
rising validation loss thereafter.
Fruits Classification (1,655 images) peaks at epoch 10 (95.8\% F1, C4) then
oscillates.
For the Korean food dataset, early stopping on validation F1 will be essential.

\end{enumerate}

% ─────────────────────────────────────────────────────────────────────────────
\section{Korean Food Dataset Proposal}
% ─────────────────────────────────────────────────────────────────────────────

\subsection{Target Categories and Classification Scheme}

Two types of Korean food items are targeted:

\textbf{Korean fruits}, directly comparable to existing datasets:
Korean apple, Jeju mandarin (\textit{gyul}), shine muscat grape,
Korean strawberry (\textit{maehyang} variety).
These share visual properties with Western counterparts but have distinct color
profiles and surface textures that make a domain-specific model preferable.

\textbf{Korean prepared foods}, with no analog in existing datasets:
Kimchi (fermented cabbage), \textit{tteok} (rice cake),
\textit{banchan} (side dishes), cooked rice (\textit{bap}).
For these categories, binary classification (fresh / rotten) is proposed,
as the intermediate ``semifresh'' state is visually ambiguous in fermented items.

\subsection{Data Collection Pipeline}

Two automated image collection scripts have been implemented, one targeting
Korean fruits and one targeting Korean prepared foods.
Both use Bing Image Search with 5--8 search queries per class per freshness
state, combining Korean-language and romanized search terms to maximize
result diversity (e.g., \textit{``fresh kimchi bright red color''} and
\textit{``moldy kimchi white fungus''} for the kimchi category).

Downloaded images are validated and processed as follows:
\begin{itemize}
    \item \textbf{Integrity check}: corrupt or unreadable files are discarded.
    \item \textbf{Minimum resolution}: $100 \times 100$ pixels.
    \item \textbf{Format}: any format accepted; converted to RGB JPEG.
    \item \textbf{Output size}: $256 \times 256$ pixels
    (model crops to $224 \times 224$ during training).
\end{itemize}

The pipeline supports incremental collection with resume capability:
if a class folder already contains images from a previous run, only the
remaining count is downloaded.
The target dataset size is 100 images per class per freshness state.
Table~\ref{tab:korean_counts} summarizes the expected dataset composition.

\begin{table}[h]
\centering
\caption{Expected composition of the Korean food dataset (target counts).}
\label{tab:korean_counts}
\begin{tabular}{lcccc}
\toprule
\textbf{Category} & \textbf{Type} & \textbf{States} & \textbf{Classes} & \textbf{Target images} \\
\midrule
Kimchi, tteok, banchan, bap & Prepared food & 2 (Fresh/Rotten)       & 8  & 800  \\
Apple, mandarin, strawberry  & Korean fruit  & 2 (Fresh/Rotten)       & 6  & 600  \\
Shine muscat, mandarin       & Korean fruit  & 3 (Fresh/Semi/Rotten)  & 6  & 600  \\
\midrule
\multicolumn{4}{l}{\textbf{Total}}                                    & \textbf{$\sim$2,000} \\
\bottomrule
\end{tabular}
\end{table}

This scale is deliberately comparable to Lemon Varieties (1,956 images),
for which the pipeline achieved 98.2\% F1 with transfer learning ---
confirming that 2,000 images is a sufficient starting point.

\subsection{Integration with Training Pipeline}

The Korean food dataset will be organized in the same
\texttt{food/state/} nested structure used throughout this study,
making it directly compatible with the existing training and evaluation
pipeline without any code modifications.
The \texttt{prepare\_datasets.py} script can be extended with a new
preparation function, and the training pipeline selects the dataset
interactively at runtime.

\end{document}
